% Compatibilidad con instalaciones MiKTeX cuyo formato base es anterior a
% los paquetes Beamer/Hyperref instalados.
\providecommand{\IfFormatAtLeastTF}[3]{#3}
\providecommand{\IfFormatAtLeastT}[2]{}
\providecommand{\IfFormatAtLeastF}[2]{#2}
\providecommand{\IfDocumentMetadataTF}[2]{#2}
\providecommand{\IfDocumentMetadataT}[1]{}
\providecommand{\IfDocumentMetadataF}[1]{#1}
\makeatletter
\@namedef{ver@epstopdf-base.sty}{9999/12/31}
\makeatother

\documentclass[aspectratio=169,11pt]{beamer}

% =============================================================================
% MACROECONOMIA II | LICENCIATURA EN ECONOMIA | CIDE
% SESION 4: CONSUMO INTERTEMPORAL Y MERCADOS DE CREDITO
% =============================================================================

% --- Idioma, tipografia y matematicas -----------------------------------------
\usepackage[utf8]{inputenc}
\IfFileExists{spanish.ldf}{
  \usepackage[spanish,es-nodecimaldot]{babel}
}{
  \usepackage[english]{babel}
}
\uselanguage{Spanish}
\languagepath{Spanish}

\usepackage{graphicx}
\usepackage{amsmath,amssymb}
\IfFileExists{booktabs.sty}{
  \usepackage{booktabs}
}{
  \providecommand{\toprule}{\hline}
  \providecommand{\midrule}{\hline}
  \providecommand{\bottomrule}{\hline}
}
\providecommand{\addlinespace}[1][0.5em]{}
\usepackage{xcolor}
\usepackage{ragged2e}
\usepackage{etoolbox}

\newtheorem{proposition}[theorem]{Proposición}
\newtheorem{remark}[theorem]{Observación}
\newtheorem{assumption}[theorem]{Supuesto}

% --- Justificacion del contenido ---------------------------------------------
\AtBeginDocument{\justifying}
\patchcmd{\itemize}{\raggedright}{\justifying}{}{}
\patchcmd{\enumerate}{\raggedright}{\justifying}{}{}
\patchcmd{\description}{\raggedright}{\justifying}{}{}
\addtobeamertemplate{block begin}{}{\justifying}
\addtobeamertemplate{block alerted begin}{}{\justifying}
\addtobeamertemplate{block example begin}{}{\justifying}
\makeatletter
\patchcmd{\beamer@columnenv}{\raggedright}{\justifying}{}{}
\patchcmd{\beamer@columncom}{\raggedright}{\justifying}{}{}
\makeatother

% --- Paleta ------------------------------------------------------------------
\definecolor{CIDENavy}{RGB}{3,25,54}
\definecolor{CIDEBlue}{RGB}{5,35,75}
\definecolor{CIDEInk}{RGB}{25,34,45}
\definecolor{CIDEGray}{RGB}{101,112,126}
\definecolor{CIDEGrayLight}{RGB}{243,244,246}

\setbeamercolor{normal text}{fg=CIDEInk,bg=white}
\setbeamercolor{structure}{fg=CIDEBlue}
\setbeamercolor{alerted text}{fg=CIDEBlue}
\setbeamercolor{title}{bg=CIDEBlue,fg=white}
\setbeamercolor{subtitle}{fg=CIDEInk}
\setbeamercolor{frametitle}{bg=CIDEBlue,fg=white}
\setbeamercolor{palette primary}{bg=CIDEBlue,fg=white}
\setbeamercolor{palette secondary}{bg=CIDEBlue,fg=white}
\setbeamercolor{palette tertiary}{bg=CIDENavy,fg=white}
\setbeamercolor{palette quaternary}{bg=CIDENavy,fg=white}
\setbeamercolor{section in head/foot}{bg=CIDENavy,fg=white}
\setbeamercolor{section in head/foot shaded}{bg=CIDENavy,fg=white!50}
\setbeamercolor{mini frame}{bg=white,fg=white}
\setbeamercolor{mini frame shaded}{bg=white!35,fg=white!35}
\setbeamercolor{upper separation line head}{bg=CIDENavy}
\setbeamercolor{lower separation line head}{bg=CIDEBlue}
\setbeamercolor{block title}{bg=CIDEBlue,fg=white}
\setbeamercolor{block body}{bg=CIDEGrayLight,fg=CIDEInk}
\setbeamercolor{block title alerted}{bg=CIDENavy,fg=white}
\setbeamercolor{block body alerted}{bg=CIDEGrayLight,fg=CIDEInk}
\setbeamercolor{caption name}{fg=CIDEBlue}

% --- Tipografia y geometria --------------------------------------------------
\setbeamerfont{author in head/foot}{size=\tiny}
\setbeamerfont{institute in head/foot}{size=\tiny}
\setbeamerfont{title in head/foot}{size=\tiny}
\setbeamersize{text margin left=0.70cm,text margin right=0.70cm}
\setlength{\leftmargini}{1.2em}
\setlength{\leftmarginii}{1.1em}
\setlength{\abovecaptionskip}{0.25em}
\setlength{\belowcaptionskip}{0pt}

\useoutertheme[subsection=false]{miniframes}
\setbeamertemplate{navigation symbols}{}
\setbeamertemplate{mini frame}[circle]
\setbeamertemplate{mini frame in current subsection}[circle]
\setbeamertemplate{section in head/foot}{\color{white}\insertsectionhead}
\setbeamertemplate{section in head/foot shaded}{\color{white!55}\insertsectionhead}
\setbeamertemplate{caption}[numbered]
\setbeamertemplate{section in toc}[circle]
\setbeamertemplate{subsection in toc}[circle]
\setbeamertemplate{itemize item}{\raisebox{0.15ex}{\tiny$\blacktriangleright$}}
\setbeamertemplate{itemize subitem}{\raisebox{0.15ex}{\tiny$\bullet$}}
\setbeamertemplate{enumerate item}[circle]

\setbeamertemplate{frametitle}{%
  \nointerlineskip
  \begin{beamercolorbox}[
    wd=\paperwidth,ht=3.05ex,dp=1.05ex,
    leftskip=0.42cm,rightskip=0.42cm
  ]{frametitle}
    \usebeamerfont{frametitle}\insertframetitle
  \end{beamercolorbox}%
}

\setbeamertemplate{footline}{%
  \leavevmode
  \vbox{%
    \offinterlineskip
    \hbox{%
      \begin{beamercolorbox}[wd=.72\paperwidth,ht=2.05ex,dp=.55ex,leftskip=.35cm]{palette secondary}
        \usebeamerfont{author in head/foot}\insertshortauthor
      \end{beamercolorbox}%
      \begin{beamercolorbox}[wd=.28\paperwidth,ht=2.05ex,dp=.55ex,rightskip=.35cm plus1fil]{palette secondary}
        \hfill\usebeamerfont{institute in head/foot}\insertshortinstitute
      \end{beamercolorbox}%
    }%
    \hbox{%
      \begin{beamercolorbox}[wd=.86\paperwidth,ht=1.85ex,dp=.50ex,leftskip=.35cm]{palette tertiary}
        \usebeamerfont{title in head/foot}\insertshorttitle
      \end{beamercolorbox}%
      \begin{beamercolorbox}[wd=.14\paperwidth,ht=1.85ex,dp=.50ex,center]{palette tertiary}
        \usebeamerfont{title in head/foot}\insertframenumber{} / \inserttotalframenumber
      \end{beamercolorbox}%
    }%
  }%
  \vskip0pt
}

\newcommand{\RutaLogoCIDE}{Figuras/General/logocide.png}
\setbeamertemplate{background}{%
  \vbox to \paperheight{%
    \vfill
    \hbox to \paperwidth{%
      \hfill
      \includegraphics[height=0.58cm,keepaspectratio]{\RutaLogoCIDE}%
      \hspace{0.18cm}%
    }%
    \vspace{0.68cm}%
  }%
}

\setbeamertemplate{title page}{%
  \vspace*{0.35cm}
  \centering
  \begin{beamercolorbox}[wd=\textwidth,sep=0.24cm,center,rounded=false]{title}
    \usebeamerfont{title}\inserttitle\par
  \end{beamercolorbox}
  \vspace{0.25cm}
  {\usebeamerfont{subtitle}\usebeamercolor[fg]{subtitle}\insertsubtitle\par}
  \vfill
  {\usebeamerfont{author}\insertauthor\par}
  \vspace{0.35cm}
  {\usebeamerfont{institute}\insertinstitute\par}
  \vfill
  {\usebeamerfont{date}\insertdate\par}
  \vspace{0.25cm}
}

% --- Comandos reutilizables --------------------------------------------------
\newcommand{\FigurasSesion}{Figuras/04_Consumo Intertemporal}
\newcommand{\figura}[2][0.65\textheight]{%
  \includegraphics[width=\linewidth,height=#1,keepaspectratio]{\FigurasSesion/#2}%
}
\newcommand{\cidehighlight}[1]{\textcolor{CIDEBlue}{\bfseries #1}}
\newcommand{\fuentefigura}[1]{%
  \vspace{0.03cm}\par{\raggedright\tiny\textit{Fuente:} #1\par}%
}
\newcommand{\notafuentes}[1]{\note{[Sources] #1}}

\hypersetup{
  colorlinks=true,
  urlcolor=blue,
  linkcolor=.,
  citecolor=.,
  filecolor=.
}

\newif\ifsectiontoc
\sectiontocfalse
\AtBeginSection[]{%
  \ifsectiontoc
    \begin{frame}{Contenido}
      \tableofcontents[currentsection,hideothersubsections]
    \end{frame}
  \fi
}

% =============================================================================
% METADATOS
% =============================================================================
\title[Macroeconomía II]{Sesión 4 \\ Consumo intertemporal y mercados de crédito}
\subtitle{Un modelo macroeconómico de dos periodos}
\author[LECO]{Arturo López}
\institute[CIDE]{Centro de Investigación y Docencia Económicas (CIDE)}
\date[Otoño 2026]{Licenciatura en Economía, Otoño 2026}

% =============================================================================
\begin{document}
% =============================================================================

\begin{frame}
  \titlepage
  \notafuentes{Williamson (2018), cap. 9; Kurlat (2020), cap. 6.}
\end{frame}

% =============================================================================
\section{Punto de partida}
% =============================================================================

\begin{frame}{Objetivos de la sesión}
  \begin{itemize}
    \item Desarrollaremos una \textbf{teoría sobre las decisiones de gasto de los hogares}.
    \item En el modelo de un periodo, el consumidor decidía cuánto consumir y cuánto trabajar. 
    \item En esta sesión introducimos una nueva decisión: consumo-ahorro.
    \item Comenzaremos el estudio de la toma de \textbf{decisiones intertemporales} y sus implicaciones para la actividad macroeconómica.
    \item Las decisiones intertemporales involucran \textit{trade-offs} a lo largo del tiempo: consumir hoy implica sacrificar consumo futuro, y viceversa.
    \item Introduciremos por primera vez un mercado financiero: el \textbf{mercado de deuda}.
    \item Verémos cómo el acceso a instrumentos de deuda permite al hogar distribuir su consumo de manera óptima entre periodos, mejorando su bienestar.
  \end{itemize}
\end{frame}

\begin{frame}{La visión keynesiana del consumo}
  Keynes (1936, citado en Kurlat, 2020, p. 103) hizo la siguiente observación sobre las decisiones de consumo:

  \vspace{0.5em}
  \textit{"La ley psicológica fundamental, en la que podemos confiar con gran seguridad tanto \textit{a priori}, por nuestro conocimiento de la naturaleza humana, como por los hechos detallados de la experiencia, es que las personas están dispuestas, por regla general y en promedio, a aumentar su consumo a medida que aumenta su ingreso, pero no en la misma medida en que aumenta el ingreso."}
  \vspace{0.5em}

  ¿Qué significa esto? ¿Es correcta esta teoría?
\end{frame}


\begin{frame}[t]{Keynes y la función de consumo}
  Traduzcamos la idea de Keynes al lenguaje matemático.
  \[
    C=c(Y),
    \qquad 0<c'(Y)<1.
  \]

  Una interpretación diferente de su argumento es que
  \[
    0<\frac{c'(Y)Y}{c(Y)}<1.
  \]
  \begin{itemize}
  \item La primera idea define una \cidehighlight{propensión marginal a consumir} positiva y menor que uno.
  \item La segunda idea propone que la elasticidad-ingreso del consumo es menor a 1.
  \end{itemize}
\end{frame}




\begin{frame}[t]{Evidencia de la función de consumo Keynesiana}
      \centering
      \figura[5.5cm]{fig_6_1_1_keynes_hogares.png}
      \fuentefigura{Kurlat (2020), figura 6.1.1; Consumer Expenditure Survey, 2014.}
\end{frame}




\begin{frame}{Implicación: sobreproducción}
  Si la función keynesiana es correcta y la elasticidad ingreso del consumo es menor que uno, entonces conforme la economía crece la razón \(C/Y\) debería ir cayendo.


  Supongamos \(C = c_0 + c_1 Y\), con \(c_0 = 100\) y \(c_1 = 0.8\):

  \vspace{0.3em}
  \begin{center}
  \begin{tabular}{cccc}
    \toprule
    \(Y\) & \(C = 100 + 0.8Y\) & \(C/Y\) & Consumo del ingreso nuevo \\
    \midrule
    500     & 500   & 1.00 & --- \\
    1{,}000  & 900   & 0.90 & de los \(500\) nuevos, se consumió el 80\% \\
    2{,}000  & 1{,}700 & 0.85 & de los \(1{,}000\) nuevos, se consumió el 80\% \\
    10{,}000 & 8{,}100 & 0.81 & de los \(8{,}000\) nuevos, se consumió el 80\% \\
    \bottomrule
  \end{tabular}
  \end{center}

  \vspace{0.5em}
  En \textbf{cada fila} se consume exactamente el 80\% del ingreso adicional. Sin embargo, \(C/Y\) cae de todos modos.

\end{frame}


% ---------- Slide 1: la conjetura ----------
\begin{frame}[t]{Implicación: sobreproducción}
\vspace{1.5cm}
  \begin{block}{La preocupación de mediados del siglo XX}
    ¿La economía va a producir cada vez más bienes que nadie quiere consumir?
    ¿Habrá desempleo masivo porque nadie desea todo lo que produciríamos si
    todos trabajaran?
  \end{block}
  
\end{frame}

% ---------- Slide 2: la evidencia ----------
\begin{frame}[t]{Evidencia agregada de la función de consumo}
  \small
  \begin{columns}[T,onlytextwidth]
    \begin{column}{0.44\textwidth}
      \begin{itemize}
        \setlength{\itemsep}{0.55em}
        \item EE.UU., 1929--2018: la elasticidad estimada es \(\approx 0.97\).
        \item Corte transversal entre países, 2011: \(\approx 0.85\).
        \item En datos agregados, el consumo ha sido una fracción casi constante del PIB (\(\approx 65\%\)). Por ejemplo, si \(c(Y)=0.65\,Y\), entonces \(\dfrac{c'(Y)\,Y}{c(Y)} = 1\).
        \item Los países \textbf{no} consumen una fracción menor de su ingreso conforme se enriquecen.
      \end{itemize}
    \end{column}
    \begin{column}{0.54\textwidth}
      \centering
      \figura[4.2cm]{fig_6_1_2_keynes_agregado.png}
      \fuentefigura{Kurlat (2020), fig.\ 6.1.2; NIPA y Feenstra et al.\ (2015).}
    \end{column}
  \end{columns}
\end{frame}


\begin{frame}{¿Por qué necesitamos una teoría intertemporal del consumo?}
  \begin{itemize}
    \setlength{\itemsep}{0.55em}
    \item El ingreso cambia a lo largo de la vida y del ciclo económico.
    \item Los hogares utilizan activos y deuda para separar el momento en que reciben ingreso del momento en que consumen.
    \item Las expectativas sobre el ingreso futuro afectan las decisiones actuales.
    \item La tasa de interés modifica el precio relativo entre consumo presente y futuro.
    \item La respuesta del consumo determina un canal central de transmisión de la política fiscal y monetaria.
  \end{itemize}
  \notafuentes{Williamson (2018), pp. 306--308; Kurlat (2020), pp. 105--106.}
\end{frame}


% =============================================================================
\section{Modelo}
% =============================================================================

\begin{frame}{Una economía de dos periodos}
  \begin{itemize}
    \setlength{\itemsep}{0.5em}
    \item Los periodos son el presente y el futuro.
    \item La economía está habitada por $N$ hogares idénticos.
    \item El hogar recibe ingresos exógenos \(y\) y \(y'\), y paga impuestos de suma fija \(t\) y \(t'\).
    \item Elige consumo \(c\), consumo futuro \(c'\) y ahorro \(s\): no hay elección de ocio.
    \item Puede comprar o emitir un bono real con rendimiento \(r\).
    \item El hogar puede prestar y pedir prestado a la misma tasa \(r\).
  \end{itemize}

  \vspace{0.12cm}
  Las primas denotan variables del periodo futuro.
  \notafuentes{Williamson (2018), pp. 308--310.}
\end{frame}



\begin{frame}{La secuencia de decisiones}
  \begin{center}
    \begin{tabular}{p{0.43\textwidth}p{0.43\textwidth}}
      \toprule
      \textbf{Periodo actual} & \textbf{Periodo futuro}\\
      \midrule
      Recibe \(y\) y paga \(t\). & Recibe \(y'\) y paga \(t'\).\\[0.15cm]
      Elige \(c\) y \(s\). & Recibe o paga \((1+r)s\).\\[0.15cm]
      Compra bonos si \(s>0\). & Consume todos sus recursos.\\[0.15cm]
      Emite bonos si \(s<0\). & Termina sin activos ni deuda.\\
      \bottomrule
    \end{tabular}
  \end{center}

  \[
    \underbrace{c+s=y-t}_{\text{periodo actual}},
    \qquad
    \underbrace{c'=y'-t'+(1+r)s}_{\text{periodo futuro}}.
  \]
  \notafuentes{Williamson (2018), ecuaciones 9-1 y 9-2.}
\end{frame}

\begin{frame}{De dos restricciones a una sola restricción intertemporal}
  \hypertarget{RP}{}
  De la restricción actual,
  \[
    s=y-t-c.
  \]
  Sustituyendo en la restricción futura,
  \[
    c'=y'-t'+(1+r)(y-t-c).
  \]
  Dividiendo entre \(1+r\) y reordenando,
  \[
    \boxed{
    c+\frac{c'}{1+r}
    =y-t+\frac{y'-t'}{1+r}
    }.
  \]
    \begin{block}{Restricción presupuestaria intertemporal}
    El valor presente del consumo debe ser igual al valor presente del ingreso disponible.
  \end{block}
  \notafuentes{Williamson (2018), pp. 310--311; Kurlat (2020), ecuación 6.2.4.}
  
  \hyperlink{app:VP}{\beamerbutton{Ver Intuición del Valor Presente}}
\end{frame}


\begin{frame}{Precio relativo: forma general para dos bienes}
  En microeconomía, de forma general teníamos que:
  $$
  p_1x_1 + p_2x_2 = m
  $$
  Dividiendo entre \(p_1\), obtenemos la recta presupuestaria en términos del bien 1:
  $$
  x_1 + \frac{p_2}{p_1}x_2 = \frac{m}{p_1}
  $$
  \(p_2/p_1\) es el precio del bien dos en unidades del bien uno. Si \(x_1\) es pan y \(x_2\) son tazas de café, con \(p_2=10\) y \(p_1=5\), tenemos que el precio relativo de una taza de café es dos unidades de pan. Es decir, a precios de mercado, una taza de café vale dos piezas de pan.
\end{frame}


\begin{frame}{Precio relativo de bienes en el tiempo (I)}
   En nuestro caso tenemos:
    \[
    c+\frac{c'}{1+r}
    =y-t+\frac{y'-t'}{1+r}.
  \]
  Normalizando el precio del consumo actual a uno, el precio relativo de una unidad de consumo futuro en unidades de consumo actual es \(\frac{1}{1+r}\).
  
  Si \(r=0.05\), el precio relativo del consumo futuro es \(\frac{1}{1.05}\approx 0.952\) unidades de consumo actual: para todo \(r>0\), el consumo futuro es más barato que el consumo actual.
\end{frame}

\begin{frame}{Precio relativo de bienes en el tiempo (II)}
  \small
  \begin{itemize}
    \setlength{\itemsep}{0.55em}
    \item $1/(1+r)$ es el precio relativo del consumo futuro en unidades de consumo actual porque renunciar hoy a 1 unidad de consumo permite obtener \(1+r\) unidades mañana.
    \item \(r\) \textbf{alto}: el consumo presente es \textit{caro} en relación al consumo futuro: el mercado ofrece muchas unidades futuras a cambio de una unidad presente.
    \item \(r\) \textbf{bajo}: el consumo presente es \textit{barato}: $1/(1+r)$ sube y el consumo futuro se vuelve relativamente más caro.
    \item La tasa real coordina la oferta de ahorro y la demanda de crédito.
  \end{itemize}

  \notafuentes{Williamson (2018), pp. 309--312; Kurlat (2020), p. 106.}
\end{frame}


\begin{frame}{Gráfica de la restricción presupuestaria intertemporal (I)}
  Para graficar la RPI,
    \[
    c+\frac{c'}{1+r}
    =y-t+\frac{y'-t'}{1+r},
  \]
  definimos la \textbf{riqueza vital} como:
  $$
  we = y - t + \frac{y' - t'}{1+r}
  $$
  Despejando,
  $$
  c' = we(1+r) - (1+r)c
  $$
  Por lo tanto,
  \begin{itemize}
    \item Intercepto en el eje c: $c = we$.
    \item Intercepto en el eje c': $c' = we(1+r)$.
    \item Pendiente: $-(1+r)$.
  \end{itemize}
\end{frame}


\begin{frame}[t]{Gráfica de la restricción presupuestaria intertemporal (II)}
  \begin{columns}[T,onlytextwidth]
    \begin{column}{0.42\textwidth}
      \small
      El punto de dotación es
      \[
        E=(y-t,\,y'-t').
      \]
      \begin{itemize}
        \item Al noroeste de \(E\): \(c<y-t\), por lo que \(s>0\).
        \item Al sureste de \(E\): \(c>y-t\), por lo que \(s<0\).
        \item La frontera reúne todas las canastas financiables mediante ahorro o deuda.
      \end{itemize}
    \end{column}
    \begin{column}{0.54\textwidth}
      \centering
      \figura[5.15cm]{fig_9_1_restriccion_intertemporal.png}
      \fuentefigura{Williamson (2018), figura 9.1.}
    \end{column}
  \end{columns}
  \notafuentes{Williamson (2018), figura 9.1 y pp. 311--313.}
\end{frame}

% =============================================================================
\section{Elección óptima}
% =============================================================================

\begin{frame}{Función de utilidad intertemporal}
  \vspace{0.5em}
  Suponemos que las preferencias del hogar están representadas por:
  $$
  U(c, c') = u(c) + \beta\, u(c')
  $$
  \vspace{-0.6cm}
  \begin{itemize}
    \item Matemáticamente idéntico al problema de dos bienes de microeconomía: aquí los ``bienes'' son \(c\) (bienes hoy) y \(c'\) (bienes en el futuro).
    \item El hogar \textbf{valora consumir hoy}, pero tambien \textbf{valora ahorrar}, porque le permite consumir más mañana.
    \item La utilidad es \textit{aditivamente separable}: la única diferencia en cómo valora cada periodo es el término \(\beta\).
    \item \(\beta\) es el \textbf{factor de descuento}, típicamente \(\beta < 1\). Representa \textit{impaciencia}: el mismo nivel de consumo da más utilidad hoy que en el futuro.
  \end{itemize}

  \notafuentes{Kurlat (2020), sección 6.2.}
\end{frame}


\begin{frame}{Supuestos sobre las preferencias}
  \begin{itemize}
    \setlength{\itemsep}{0.5em}
    \item Suponemos que \(c, c'\) son bienes normales.
    \item Además \(U_c>0\), \(U_{c'}>0\) (más consumo siempre es mejor) y
      \(U_{cc}<0\), \(U_{c'c'}<0\) (utilidad marginal decreciente en cada
      periodo).
    \item Bajo separabilidad aditiva, esto implica curvas de indiferencia
      estrictamente convexas: canastas equilibradas son preferidas a los extremos.
    \item Esa convexidad hace que los hogares deseen \textbf{suavizar su consumo en el tiempo}: preferir niveles de consumo similares en ambos periodos, en vez de mucho en uno y poco en otro.
  \end{itemize}
\end{frame}



\begin{frame}{Suavización del consumo}
  La utilidad marginal decreciente hace costoso concentrar el consumo en un
  solo periodo. Esto implica que el hogar busque suavizar su consumo.

  \begin{itemize}
    \setlength{\itemsep}{0.55em}
    \item El consumo se suaviza \textbf{independientemente} de la trayectoria del ingreso: si hay brincos en el ingreso, el ahorro (o la deuda) le permite al hogar suavizar su consumo.
    \item Si el ingreso es alto hoy y bajo mañana, el hogar \textbf{ahorra}.
    \item Si el ingreso es bajo hoy y alto mañana, el hogar \textbf{se endeuda}.
  \end{itemize}

  \vspace{0.4em}
  \cidehighlight{El consumo responde al \textbf{valor presente} del ingreso
  total en el tiempo, no únicamente al ingreso actual.}
\end{frame}



\begin{frame}{El problema del hogar}
  El hogar resuelve
  \[
    \max_{c,c'>0}
    \left\{u(c)+\beta u(c')\right\}
  \]
  sujeto a
  \[
    c+\frac{c'}{1+r}
    \leq
    y-t+\frac{y'-t'}{1+r}.
  \]

  El lagrangiano es
  \[
    \mathcal L
    =u(c)+\beta u(c')
    +\lambda\left[
      y-t+\frac{y'-t'}{1+r}
      -c-\frac{c'}{1+r}
    \right].
  \]
  \notafuentes{Kurlat (2020), ecuación 6.2.5.}
\end{frame}

\begin{frame}{Condiciones de primer orden}
  Para una solución interior,
  \begin{align*}
    u'(c)-\lambda&=0,\\
    \beta u'(c')-\frac{\lambda}{1+r}&=0,\\
    y-t+\frac{y'-t'}{1+r}-c-\frac{c'}{1+r}&=0.
  \end{align*}

  Eliminando \(\lambda\), obtenemos
  \[
    \boxed{u'(c)=\beta(1+r)u'(c')}.
  \]

  Esta condición de optimalidad intertemporal se conoce como la \cidehighlight{ecuación de Euler}.
  \notafuentes{Kurlat (2020), ecuaciones 6.2.6--6.2.8.}
\end{frame}

\begin{frame}{Interpretación de la ecuación de Euler}

  \begin{itemize}
    \item Intuición: consumir $\$1$ hoy vs ahorrar $\$1$ y consumirlo después.
    \begin{itemize}
      \item Consume $\$1$ hoy: la utilidad incrementa por $u'(c)$.
      \item Ahorra $\$1$: recibe $(1+r)$ mañana, la utilidad futura incrementa $\beta u'(c')$.
    \end{itemize}
    \item En el óptimo, el hogar es indiferente entre ambas opciones.
  \end{itemize}
  \[
    \underbrace{u'(c)}_{\substack{\text{costo marginal}\\\text{en utilidad}}}
    =
    \underbrace{\beta(1+r)u'(c')}_{\substack{\text{beneficio marginal}\\\text{descontado}}}.
  \]

  \begin{block}{Condición de tangencia}
    \small
    La tasa marginal de sustitución intertemporal es igual al precio relativo: \(u'(c)/[\beta u'(c')]=1+r\).
  \end{block}
  \notafuentes{Kurlat (2020), pp. 107--108.}
\end{frame}



\begin{frame}[t]{Prestamistas y deudores}
  \begin{columns}[T,onlytextwidth]
    \begin{column}{0.49\textwidth}
      \centering
      \figura[4.4cm]{fig_9_3_prestamista.png}
      \fuentefigura{Williamson (2018), figura 9.3.}
      \vspace{0.08cm}
      \small\(c^*<y-t\Rightarrow s^*>0\): el hogar es prestamista.
    \end{column}
    \begin{column}{0.49\textwidth}
      \centering
      \figura[4.4cm]{fig_9_4_deudor.png}
      \fuentefigura{Williamson (2018), figura 9.4.}
      \vspace{0.08cm}
      \small\(c^*>y-t\Rightarrow s^*<0\): el hogar es deudor.
    \end{column}
  \end{columns}
  \notafuentes{Williamson (2018), figuras 9.3--9.4 y pp. 315--317.}
\end{frame}


\begin{frame}{Decisiones de consumo}
      Decisiones de consumo ante dos trayectorias de ingreso diferentes
      \centering
      \figura[5.5cm]{fig_6_2_2_dos_ejemplos.png}
      \fuentefigura{Kurlat (2018), figura 6.2.2}
\end{frame}


\begin{frame}{Ejemplo: preferencias CRRA}
  Supongamos que
  \[
    u(c)=\frac{c^{1-\sigma}-1}{1-\sigma},
    \qquad \sigma>0,
  \]
  Entonces, la ecuación de Euler es:
  \[
    c^{-\sigma}=\beta(1+r)(c')^{-\sigma},
  \]
  Equivalentemente,
  \[
    \frac{c'}{c}=[\beta(1+r)]^{1/\sigma}.
  \]

  \begin{itemize}
    \item \(1/\sigma\) es la elasticidad de sustitución intertemporal.
    \item Un \(\sigma\) alto implica poca disposición a sustituir consumo entre periodos.
  \end{itemize}
  \notafuentes{Kurlat (2020), pp. 110--111.}
\end{frame}

\begin{frame}{La regla de decisión bajo preferencias CRRA}
  Sustituyendo
  \(
    c'=[\beta(1+r)]^{1/\sigma}c
  \)
  en la restricción intertemporal,
  \[
    \boxed{
    c(r,y,y',t,t')
    =
    \frac{
      y-t+\dfrac{y'-t'}{1+r}
    }{
      1+\beta^{1/\sigma}(1+r)^{1/\sigma-1}
    }
    }.
  \]

  \begin{itemize}
    \item La demanda de consumo presente depende del ingreso disponible en ambos periodos y de la tasa de interés real.
    \item El numerador es la riqueza intertemporal: el consumo presente es propocional a esta.
    \item El denominador distribuye esa riqueza entre ambos periodos.
  \end{itemize}
  \notafuentes{Kurlat (2020), ecuación 6.2.9.}
\end{frame}


% =============================================================================
\section{Ingreso}
% =============================================================================



\begin{frame}{Un aumento del ingreso actual}
  \begin{block}{}
    ¿Cómo afecta a la decisión intertemporal de consumo-ahorro del hogar un aumento en el ingreso actual?
  \end{block}

  Considere \(\Delta y>0\), manteniendo \(y'\), \(t\), \(t'\) y \(r\) constantes.

  \begin{itemize}
    \setlength{\itemsep}{0.52em}
    \item La riqueza intertemporal aumenta exactamente en \(\Delta y\).
    \item La restricción presupuestaria se desplaza paralelamente hacia afuera (efecto ingreso).
    \item Como ambos consumos son bienes normales, \(c\uparrow\) y \(c'\uparrow\).
    \item Como el hogar distribuye el incremento entre ambos periodos,
      \[
        0<\Delta c<\Delta y
        \quad\Longrightarrow\quad
        \Delta s=\Delta y-\Delta c>0.
      \]
  \end{itemize}
  \notafuentes{Williamson (2018), pp. 317--319.}
\end{frame}

\begin{frame}[t]{El hogar distribuye un aumento transitorio del ingreso}
  \begin{columns}[T,onlytextwidth]
    \begin{column}{0.40\textwidth}
      \small
      El ingreso actual aumenta de \(y_1\) a \(y_2\):
      \begin{itemize}
        \item la pendiente no cambia;
        \item \(c\) aumenta de \(c_1\) a \(c_2\);
        \item \(c'\) también aumenta;
        \item parte del incremento se ahorra para suavizar el consumo.
      \end{itemize}
    \end{column}
    \begin{column}{0.56\textwidth}
      \centering
      \figura[5.1cm]{fig_9_5_ingreso_actual.png}
      \fuentefigura{Williamson (2018), figura 9.5.}
    \end{column}
  \end{columns}
  \notafuentes{Williamson (2018), figura 9.5 y pp. 317--319.}
\end{frame}

\begin{frame}{Un aumento del ingreso futuro}
  Considere \(\Delta y'>0\), manteniendo las demás variables constantes.
  \[
    \Delta\omega=\frac{\Delta y'}{1+r}>0.
  \]

  \begin{itemize}
    \setlength{\itemsep}{0.55em}
    \item Aumentan el consumo actual y el futuro.
    \item El hogar puede elevar \(c\) antes de recibir el ingreso mediante endeudamiento.
    \item Por tanto, el ahorro actual disminuye:
      \[
        \Delta s=-\Delta c<0.
      \]
  \end{itemize}

  \cidehighlight{La expectativa de mayor ingreso futuro modifica decisiones actuales}.
  \notafuentes{Williamson (2018), pp. 321--323.}
\end{frame}

\begin{frame}[t]{Las expectativas sobre ingreso futuro importan hoy}
  \begin{columns}[T,onlytextwidth]
    \begin{column}{0.40\textwidth}
      \small
      Cuando \(y'\) aumenta:
      \begin{itemize}
        \item la riqueza aumenta;
        \item la restricción se desplaza hacia afuera;
        \item \(c\) y \(c'\) aumentan;
        \item el hogar ahorra menos o se endeuda más.
      \end{itemize}
    \end{column}
    \begin{column}{0.56\textwidth}
      \centering
      \figura[5.0cm]{fig_9_8_ingreso_futuro.png}
      \fuentefigura{Williamson (2018), figura 9.8.}
    \end{column}
  \end{columns}
  \notafuentes{Williamson (2018), figura 9.8 y pp. 321--323.}
\end{frame}



\begin{frame}{Cambios transitorios y permanentes en el ingreso}
  \begin{itemize}
    \setlength{\itemsep}{0.6em}
    \item Hasta ahora hemos analizado el efecto de cambios en el ingreso de un periodo determinado.
    \item Sin embargo, cuando un hogar recibe un cambio en su ingreso actual, es importante para su decisión de consumo si ese cambio es \textbf{temporal} o \textbf{permanente}.
    \item La distinción entre el efecto de cambios transitorios y
      permanentes en el ingreso fue propuesta por Milton Friedman (1957) en su \textbf{hipótesis del ingreso permanente} (HIP).
    \item \textbf{Keynes}: el principal determinante del consumo es el \textbf{ingreso corriente}.
    \item \textbf{HIP}: el consumo depende del \textbf{ingreso corriente y del ingreso esperado en el futuro}.
  \end{itemize}
\end{frame}

\begin{frame}{Friedman (1957) “A Theory of the Consumption Function”}
  \centering
  \figura[6.0cm]{Friedman.png}
\end{frame}

\begin{frame}{Un ejemplo}
  Supongamos que una persona recibe un incremento de \$1{,}000 en su ingreso.
  ¿Responde igual si es temporal que si es permanente?

  \vspace{0.5em}
  \begin{columns}[T,onlytextwidth]
    \begin{column}{0.48\textwidth}
      \textbf{Incremento transitorio}\\[0.3em]
      Un bono único de \$1{,}000 este año.
      \begin{itemize}
        \item El consumo aumenta \textbf{poco}.
        \item El resto se \textbf{ahorra}, para suavizar el consumo en el
          tiempo.
      \end{itemize}
    \end{column}
    \begin{column}{0.48\textwidth}
      \textbf{Incremento permanente}\\[0.3em]
      Un aumento de \$1{,}000 al salario base, que se espera continúe
      indefinidamente.
      \begin{itemize}
        \item El consumo aumenta \textbf{mucho más}.
      \end{itemize}
    \end{column}
  \end{columns}

  \vspace{0.6em}
  \cidehighlight{La PMgC de un ingreso permanente es mayor que la PMgC de un
  ingreso transitorio.}

  \notafuentes{Kurlat (2020), sección 6.3.}
\end{frame}



\begin{frame}[t]{Análisis gráfico de la HIP}
  \begin{columns}[T,onlytextwidth]
    \begin{column}{0.5\textwidth}
      \footnotesize
      \begin{itemize}
        \item Considere un aumento temporal en el ingreso corriente del hogar: $\Delta y = y_2-y_1$, de H a L.
        \item Para suavizar su consumo, el hogar incrementa su consumo corriente de $c_1$ a $c_2$ y ahorra el resto.
        \item Ahora supongamos que el hogar espera que el incremento sea permanente. Entonces, $y'_2-y'_1 = y_2-y_1$, de L a M.
        \item El hogar incrementa su consumo corriente a $c_3$.
        \item Noten que $c_3 > c_2$. Sin embargo, ¡el incremento en el ingreso corriente es el mismo que antes! $y_2-y_1$.
        \item Lo único que cambió es que ahora el hogar espera que el incremento en el ingreso actual sea permanente.
      \end{itemize}
    \end{column}
    \begin{column}{0.48\textwidth}
      \centering
      \figura[5.05cm]{fig_9_9_temporal_permanente.png}
      \fuentefigura{Williamson (2018), figura 9.9.}
    \end{column}
  \end{columns}
  \notafuentes{Williamson (2018), figura 9.9 y pp. 323--325.}
\end{frame}



\begin{frame}{Ejemplo analítico de la HIP (I)}
   Bajo preferencias CRRA, la propensión marginal a consumir (PMgC) del ingreso transitorio es:
  \[
    \frac{\partial c}{\partial y}=\frac{1}{1+\beta^{1/\sigma}(1+r)^{1/\sigma-1}} < 1,
    \quad \text{o} \quad 
    \frac{\partial c}{\partial y'}=\frac{1}{1+r} \cdot \frac{1}{1+\beta^{1/\sigma}(1+r)^{1/\sigma-1}} <1.
  \]
  Si \(y\) y \(y'\) aumentan conjuntamente en \(\Delta\),
  \[
    \frac{dc}{d\Delta}
    =\frac{1}{1+\beta^{1/\sigma}(1+r)^{1/\sigma-1}} \left(1+\frac{1}{1+r}\right)
    >\frac{\partial c}{\partial y}.
  \]

  La propensión marginal a consumir es mayor cuando el cambio en el ingreso es permanente.
  \notafuentes{Kurlat (2020), ecuaciones 6.2.10--6.2.11.}
\end{frame}


\begin{frame}{Ejemplo analítico de la HIP (II)}
   En particular, si $\beta (1 + r) = 1$:
  \[
    \frac{dc}{d\Delta}
    =\frac{1}{1+\beta^{1/\sigma}(1+r)^{1/\sigma-1}} \left(1+\frac{1}{1+r}\right)
    = 1
  \]
  \cidehighlight{En otras palabras, cuando el hogar percibe un incremente permanente en su ingreso, este incrementa su consumo en la misma magnitud del incremento en su ingreso}.
\end{frame}


\begin{frame}{Comparación respecto a la función Keynesiana}
  Bajo CRRA, vimos que la PMgC respecto a choques de ingreso transitorio corriente es
    \[
    \frac{\partial c}{\partial y}=\frac{1}{1+\beta^{1/\sigma}(1+r)^{1/\sigma-1}}
  \]
  En contraste, \textbf{la función Keynesiana}: $C = c_0 + c_1(y-t)$, donde $PMgC = c_1$ dado exógenamente.

  En nuestro modelo, la PMgC es endógena y depende de las preferencias $(\beta, \sigma)$ y de los precios $r$.

  La función Keynesiana es una \textbf{función no microfundada}, sujeta a la crítica de Lucas (que veremos en próximas clases): cambios en $r$ resultan en cambios en la PMgC.
\end{frame}


\begin{frame}{Resultados principales de la HIP}
  \begin{itemize}
    \item Los cambios \textbf{transitorios} del ingreso generan respuestas pequeñas del consumo: el ahorro absorbe buena parte de las fluctuaciones transitorias.
    \item Los cambios \textbf{permanentes} generan respuestas mucho mayores, porque desplazan el valor presente del ingreso total.
    \item La distinción relevante no es "ingreso alto vs. bajo", sino \textbf{choques permanentes vs. transitorios} en el ingreso.
  \end{itemize}

  \cidehighlight{Este mecanismo reconcilia la relación consumo-ingreso observada entre hogares con la relación consumo-ingreso casi proporcional en los datos agregados y de largo plazo.}
\end{frame}


\begin{frame}[t]{La HIP y los datos micro}
  En un corte transversal de hogares, el ingreso \textit{total} observado mezcla ingreso permanente con choques transitorios (buenos y malos años).

  \vspace{0.3em}
  \centering
  \figura[5.0cm]{fig_6_1_1_keynes_hogares.png}
  \fuentefigura{Kurlat (2020), figura 6.1.1; Consumer Expenditure Survey, 2014.}

  \vspace{0.3em}
  \raggedright
  Quienes tuvieron un \textit{buen año} consumen relativamente poco de ese ingreso extra, porque saben que es en parte transitorio.

  \notafuentes{Kurlat (2020), figura 6.1.1.}
\end{frame}


\begin{frame}[t]{La HIP y los datos macro}
  En datos agregados y/o de largo plazo, los choques transitorios se \textbf{cancelan}: lo que queda es, esencialmente, ingreso permanente.

  \vspace{0.3em}
  \centering
  \figura[5.0cm]{fig_6_1_2_keynes_agregado.png}
  \fuentefigura{Kurlat (2020), figura 6.1.2; NIPA y Feenstra et al. (2015).}

  \notafuentes{Kurlat (2020), figura 6.1.2.}
\end{frame}



\begin{frame}[t]{Suavización de consumo en US}
  \begin{columns}[T,onlytextwidth]
    \begin{column}{0.48\textwidth}
      \centering
      \figura[5.2cm]{fig_9_6_durables.png}
      \fuentefigura{Williamson (2018), figura 9.6.}
    \end{column}
      \begin{column}{0.48\textwidth}
      \centering
      \figura[5.2cm]{fig_9_7_no_durables.png}
      \fuentefigura{Williamson (2018), figura 9.7.}
    \end{column}
    \end{columns}
\end{frame}

\begin{frame}{Suavización de consumo en México}
  \centering
      \figura[6.8cm]{consumo-mx.png}
\end{frame}



\begin{frame}{Más allá de la HIP}
  \begin{itemize}
    \item La HIP propone que la propensión marginal a consumir varía en función al tipo de ingreso que recibe el choque: choques en el ingreso permanente generan una mayor propensión marginal a consumir, cercana a uno, mientras choques en el ingreso transitorio genera una PMgC muy pequeña ($PMgC \approx 0.05$ en calibraciones típicas).
    \item Sin embargo, esta no es la única fuente potencial de variación de la propensión marginal a consumir.
    \item Por ejemplo, la PMgC puede variar por la cantidad de reservas de activos líquidos que tienen los hogares o de acuerdo a su raza (Ganong et al., 2020).
    \item A su vez, estimaciones empíricas de la PMgC del ingreso transitorio son mayores a las predicciones de la HIP: un punto medio entre Friedman y Keynes ($\approx 0.30$).
  \end{itemize}
\end{frame}


\begin{frame}{Fallas de la HIP: Exceso de Sensibilidad del Gasto al Ingreso}
\small

\begin{itemize}

    \item \textbf{Modelos racionales con restricciones de liquidez}
    \begin{itemize}
        \item Agentes \textit{forward-looking} tipo HIP, pero cerca de su límite de endeudamiento.
        \item No pueden suavizar consumo \textbf{pidiendo prestado} antes de recibir el ingreso.
        \item Predicción: el gasto futuro aumenta cuando se experimenta un incremento anticipado en el ingreso futuro.
    \end{itemize}


    \item \textbf{Modelos conductuales}
    \begin{itemize}
        \item Agentes con sesgo al presente (\textit{present bias}) o con \textit{miopía}.
        \item No planifican correctamente el futuro.
        \item Predicción: el gasto también aumenta cuando llega el ingreso anticipado.
    \end{itemize}

    \item \textbf{Ante un aumento predecible del ingreso, ambas teorías son observacionalmente equivalentes: el gasto responde excesivamente porque los hogares son conductuales o porque enfrentan restricciones de liquidez.}

\end{itemize}

\end{frame}


\begin{frame}{Ganong et al., 2020}
    \centering
      \figura[6.7cm]{Ganong et al.png}
\end{frame}

\begin{frame}{PMgC en función de las reservas de activos líquidos}
    \centering
      \figura[6.7cm]{Ganong_fig.png}
\end{frame}


\begin{frame}{Sesgos conductuales}
  \begin{itemize}
    \item Hasta ahora hemos asumido que los hogares son racionales. Esto implica que son \textit{forward-looking}: contemplan el ingreso presente y futuro en su toma de decisiones de consumo.
    \item Sin embargo, sesgos conductuales podrían determinar las decisiones de gasto de los hogares.
    \item Por ejemplo, hogares con \textit{sesgo al presente} o con \textit{miopía} le darían poco peso al ingreso futuro en sus decisiones de consumo presente.
    \item Lo anterior genera conductas \textit{hand-to-mouth} que implican PMgC cercanas a uno para choques transitorios de ingreso. 
    \item Esto resulta en multiplicadores fiscales más cercanos a los Keynesianos que a los que predicen la HIP
  \end{itemize}
\end{frame}

\begin{frame}{Ganong y Noel (2019)}
  \centering
  \figura[6.7cm]{Ganong_UI.png}
\end{frame}

\begin{frame}{Ausencia de suavización del consumo}
  \centering
  \figura[6.7cm]{Ganong_UI_fig.png}
\end{frame}

\begin{frame}{Resultados principales de Ganong y Noel (2019)}
  \begin{itemize}
    \item Una PMgC agregada estimada de aproximadamente $0.27$.
    \item Encuentran que un modelo de agentes heterogéneos ajusta los datos mucho mejor que un modelo de agente racional.
    \item Estiman que aproximadamente el 25\% de agentes tienen una conducta \textit{hand-to-mouth}.
    \item El resto son una combinación de agentes racionales y agentes con restricciones de liquidez.
  \end{itemize}
\end{frame}


\begin{frame}{Reparando a la HIP: literatura HANK}
  \begin{itemize}
    \item La forma de mejorar los modelos de consumo intertemporal es incluyendo esta \textbf{heterogeneidad en la PMgC}.
    \item Lo anterior se puede hacer al incluir hogares con distintas características: hogares racionales, hogares con sesgos conductuales, hogares con restricciones de liquidez.
    \item Es decir, \textbf{ir más allá del agente representativo}.
  \end{itemize}
\end{frame}

\begin{frame}{Macro y heterogeneidad}
  La \textbf{idea central}: \textit{la macro importa para la distribución y la distribución afecta a la macro}.
  \vspace{0.6cm}
  \begin{columns}
    \begin{column}{0.48\textwidth}
        \centering
        \figura[5cm]{MP_HANK.png}
    \end{column}
        \begin{column}{0.48\textwidth}
        \centering
        \figura[5cm]{MP_HANK_fig.png}
    \end{column}
  \end{columns}
\end{frame}


% =============================================================================
\section{Tasa real}
% =============================================================================

\begin{frame}{Efectos de la tasa de interés real}

  \begin{itemize}
    \item Hasta ahora, nos hemos enfocado en la \textbf{respuesta del consumo ante cambios en la trayectoria del ingreso (PMgC)}.
    \item Este mecanismo será clave para estudiar los efectos y mecanismos de propagación de la \textbf{política fiscal} y \textbf{choques de productividad}.
    \item Ahora, nos enfocaremos en la \textbf{respuesta del consumo ante cambios en la tasa de interés real}.
    \item Este mecanismo será clave para estudiar los efectos y mecanismos de propagación de la \textbf{política monetaria}.
  \end{itemize}

\end{frame}


\begin{frame}{Un aumento en la tasa de interés real}
  Si \(r_2>r_1\), el precio del consumo futuro en unidades de consumo presente cae:
  \[
    \frac{1}{1+r_2}<\frac{1}{1+r_1}.
  \]

  \begin{itemize}
    \setlength{\itemsep}{0.52em}
    \item La restricción gira alrededor del punto de dotación \((y-t,y'-t')\).
    \item Su pendiente pasa de \(-(1+r_1)\) a \(-(1+r_2)\).
    \item El cambio genera un efecto sustitución y un efecto ingreso (como cambios en $w$ en el problema de ocio-consumo).
    \item El efecto ingreso depende de la posición financiera inicial del hogar.
  \end{itemize}
  \notafuentes{Williamson (2018), pp. 327--329; Kurlat (2020), pp. 109--110.}
\end{frame}

\begin{frame}[t]{La restricción gira alrededor de la dotación}
  \begin{columns}[T,onlytextwidth]
    \begin{column}{0.42\textwidth}
      \small
      \begin{itemize}
        \item El hogar siempre puede consumir su ingreso disponible en cada fecha. Por ello, el punto de dotación sigue siendo factible para cualquier \(r\).
        \item Una tasa mayor recompensa más el ahorro y encarece el endeudamiento.
      \end{itemize}

      \vspace{0.15cm}
    
    \end{column}
    \begin{column}{0.54\textwidth}
      \centering
      \figura[5.2cm]{fig_9_12_tasa_interes.png}
      \fuentefigura{Williamson (2018), figura 9.12.}
    \end{column}
  \end{columns}
  \notafuentes{Williamson (2018), figura 9.12 y pp. 327--329.}
\end{frame}

\begin{frame}{Efectos sustitución e ingreso}
  \begin{itemize}
    \setlength{\itemsep}{0.55em}
    \item \textbf{Efecto sustitución:} al aumentar \(r\), el consumo actual se vuelve relativamente más caro.
      \[
        c\downarrow,
        \qquad c'\uparrow,
        \qquad s\uparrow.
      \]
    \item \textbf{Efecto ingreso para un prestamista:} recibe un mayor rendimiento sobre su ahorro; su riqueza intertemporal aumenta.
    \item \textbf{Efecto ingreso para un deudor:} aumenta el costo de su deuda; su riqueza intertemporal disminuye.
  \end{itemize}

  \begin{alertblock}{Resultado general}
    La respuesta total del consumo y del ahorro a la tasa real no puede determinarse sin conocer preferencias y posición financiera.
  \end{alertblock}
  \notafuentes{Williamson (2018), pp. 328--332; Kurlat (2020), pp. 109--110.}
\end{frame}

\begin{frame}[t]{Una tasa mayor para un ahorrador}
  \begin{columns}[T,onlytextwidth]
    \begin{column}{0.39\textwidth}
      \small
      \begin{itemize}
        \item Efecto sustitución: \(c\downarrow\), \(c'\uparrow\).
        \item Efecto ingreso: \(c\uparrow\), \(c'\uparrow\).
        \item \(c'\) aumenta inequívocamente.
        \item El efecto sobre \(c\) y \(s\) es ambiguo.
        \item Si el efecto ingreso domina al efecto sustitución, entonces $c$ aumenta (por suavización del consumo).
        \item Si el efecto sustitución domina al efecto ingreso, el ahorro aumenta de modo que $c$ cae.
      \end{itemize}
    \end{column}
    \begin{column}{0.57\textwidth}
      \centering
      \figura[5.5cm]{fig_9_13_tasa_prestamista.png}
      \fuentefigura{Williamson (2018), figura 9.13.}
    \end{column}
  \end{columns}
  \notafuentes{Williamson (2018), figura 9.13 y pp. 329--331.}
\end{frame}


\begin{frame}{Ejemplo: Efectos ingreso y sustitución para un ahorrador}

\begin{itemize}
\small
    \item Christine es inicialmente \textbf{ahorradora}:
    \[
    y = y' = \$40{,}000,
    \qquad r_1 = 5\%,
    \qquad s = 0.30\cdot y = \$12{,}000.
    \]

    Por lo tanto,
    \[
    c = \$28{,}000,
    \qquad
    c' = 40{,}000 + (1.05)(12{,}000) = \$52{,}600.
    \]

    \item Supongamos ahora que la tasa de interés real aumenta a
    \[
    r_2 = 10\%.
    \]

    \item \textbf{Efecto sustitución:} el consumo futuro se vuelve relativamente más barato:
    \[
    c' = 40{,}000 + (1.10)(12{,}000) = \$53{,}200.
    \]
    Esto incentiva a Christine a sustituir consumo presente por consumo futuro.

    \item \textbf{Efecto ingreso:} como Christine es ahorradora, el aumento en $r$ incrementa su riqueza. Para alcanzar el consumo futuro original:
    \[
    s = \frac{52{,}600-40{,}000}{1.10} = \$11{,}454,
    \qquad
    c = 40{,}000-11{,}454 = \$28{,}546.
    \]
    Esto incentiva un mayor consumo presente.

\end{itemize}

\end{frame}

\begin{frame}[t]{Una tasa mayor para un deudor}
  \begin{columns}[T,onlytextwidth]
    \begin{column}{0.39\textwidth}
      \small
      \begin{itemize}
        \item Efecto sustitución: \(c\downarrow\), \(c'\uparrow\).
        \item Efecto ingreso: \(c\downarrow\), \(c'\downarrow\).
        \item \(c\) disminuye inequívocamente.
        \item El efecto sobre \(c'\) es ambiguo.
        \item El endeudamiento disminuye: \(s\uparrow\).
      \end{itemize}
    \end{column}
    \begin{column}{0.57\textwidth}
      \centering
      \figura[5.5cm]{fig_9_14_tasa_deudor.png}
      \fuentefigura{Williamson (2018), figura 9.14.}
    \end{column}
  \end{columns}
  \notafuentes{Williamson (2018), figura 9.14 y pp. 330--332.}
\end{frame}


\begin{frame}{Ejemplo: Efectos ingreso y sustitución para un deudor}
\small
\begin{itemize}

    \item Christopher es inicialmente \textbf{deudor}:
    \[
    y = y' = \$40{,}000,
    \qquad r_1 = 5\%,
    \qquad s = -\$20{,}000.
    \]

    Por lo tanto,
    \[
    c = \$60{,}000,
    \qquad
    c' = 40{,}000 - (1.05)(20{,}000) = \$19{,}000.
    \]

    \item Supongamos ahora que la tasa de interés real aumenta a
    \[
    r_2 = 10\%.
    \]

    \item \textbf{Efecto sustitución:} el costo de la deuda aumenta. Esto incentiva a Christopher a endeudarse menos (sustituir consumo presente por consumo futuro).

    \item \textbf{Efecto ingreso:} como Christopher es deudor, el aumento en $r$ reduce su riqueza. Para mantener el consumo futuro original:
    \[
    s = \frac{40{,}000-19{,}000}{1.10} = -\$19{,}091,
    \qquad
    c = 40{,}000+19{,}091 = \$59{,}091.
    \]

\end{itemize}

\end{frame}



\begin{frame}{Efecto sustitución e ingreso según posición financiera}
  \centering
  \figura[5.2cm]{EI_ES.png}
\end{frame}






\begin{frame}{Efectos individuales de la tasa de interés real}

  Ante un aumento en la tasa de interés real, hay un \textbf{efecto de sustitución intertemporal} para ambas posiciones financieras. El efecto ingreso varía de acuerdo a la posición:

  \begin{itemize}

    \item \textbf{Efecto sustitución:} opera en la misma dirección para
    ahorradores y deudores:
    \[
      r \uparrow
      \quad\Rightarrow\quad
      \frac{1}{1+r}\downarrow
      \quad\Rightarrow\quad
      c\downarrow,\; c'\uparrow
      \quad\Rightarrow\quad
      s\uparrow.
    \]

    \item Una mayor tasa de interés real reduce el precio relativo del consumo
    futuro en términos de consumo presente, induciendo a sustituir consumo
    presente por consumo futuro.

    \item \textbf{Efecto ingreso:} depende de la posición financiera inicial:
    \begin{itemize}
      \item \textbf{Ahorradores:}
      \[
        r\uparrow \quad\Rightarrow\quad
        \text{riqueza}\uparrow \quad\Rightarrow\quad c\uparrow.
      \]

      \item \textbf{Deudores:}
      \[
        r\uparrow \quad\Rightarrow\quad
        \text{riqueza}\downarrow \quad\Rightarrow\quad c\downarrow.
      \]
    \end{itemize}

  \end{itemize}

\end{frame}


\begin{frame}{Efectos agregados de la tasa de interés real}

  En macroeconomía, nos interesa principalmente la respuesta del
  \textbf{consumo agregado} ante cambios en la tasa de interés real.

  \begin{itemize}

    \item Para los \textbf{deudores}, ambos efectos reducen el consumo presente. Para los \textbf{ahorradores}, los efectos ingreso y sustitución operan en direcciones opuestas sobre el consumo presente.

    \item En el agregado, los efectos ingreso negativos sobre el consumo de los deudores tienden a compensar los efectos ingreso positivos sobre el consumo de los ahorradores, dejando principalmente el efecto sustitución.

    \item  Sin embargo, esta compensación no tiene por qué ser exacta: \textbf{no existe una garantía teórica de que el consumo agregado
    disminuya cuando aumenta la tasa de interés real}.

  \end{itemize}

\end{frame}



% =============================================================================
\section{Mercado de crédito}
% =============================================================================

\begin{frame}{Equilibrio general dinámico}

  \begin{itemize}

    \item Ahora incorporamos al gobierno para construir nuestro primer modelo dinámico de equilibrio general (DGE).
    
    \item El modelo es una economía sin producción, compuesta por $N$ hogares y gobierno, que interactúan en dos mercados: el mercado de crédito y el mercado de bienes y servicios.
    
    \item En el mercado de crédito, hogares y gobierno pueden prestar y pedir prestado a través de un bono.

    \item Nos enfocaremos en estudiar cómo la política fiscal afecta las decisiones intertemporales de consumo y ahorro de los hogares.

    \item En particular, analizaremos un resultado central y ampliamente discutido en macroeconomía y economía pública: la \textbf{equivalencia ricardiana}.

  \end{itemize}

\end{frame}



\begin{frame}{Restricción presupuestaria intertemporal del gobierno}
  En cada periodo, el gobierno está sujeto a las siguientes restricciones presupuestarias:
  \[
    G=T+B,
    \qquad
    G'+(1+r)B=T',
  \]
  donde $T=Nt$ y $T'=Nt'$ es la recaudación total del gobierno y \(B>0\) es deuda pública emitida en el periodo presente.

  Eliminando \(B\), obtenemos la \textbf{restricción presupuestaria intertemporal del gobierno}:
  \[
    \boxed{
    G+\frac{G'}{1+r}
    =T+\frac{T'}{1+r}
    }
  \]
  que muestra que \textbf{el valor presente del gasto público debe ser igual al valor presente de los impuestos recaudados a lo largo del tiempo}.

  Es decir, el gobierno debe eventualmente pagar toda su deuda a través de cobrar impuestos a los hogares.
  \notafuentes{Williamson (2018), ecuaciones 9-17--9-19.}
\end{frame}

\begin{frame}{Equilibrio competitivo de la economía de dos periodos}
  Dadas las variables exógenas, un equilibrio competitivo consiste en una asignación \((c_i^*,c_i'^*)_{i=1}^N\), una posición financiera $(s_i^*)_{i=1}^N$ y una tasa de interés real \(r^*\) tales que:
  \begin{enumerate}
    \setlength{\itemsep}{0.55em}
    \item para cada hogar $i=1, \dots, N$, \((c_i^*,c_i'^*)\) resuelve su problema de maximización de utilidad, tomando \(r^*\) como dado;
    \item el gobierno satisface su restricción presupuestaria intertemporal;
    \item el mercado de crédito y el mercado de bienes se vacían.
  \end{enumerate}

  \begin{block}{Variable de ajuste}
    La tasa de interés real se determina endógenamente para hacer compatibles los planes de ahorro y endeudamiento.
  \end{block}
  \notafuentes{Williamson (2018), pp. 335--336.}
\end{frame}

\begin{frame}{Vaciado del mercado de crédito}
  El ahorro privado agregado es
  \[
    S^p(r)=\sum_{i=1}^{N}s_i(r).
  \]
  La condición de equilibrio es
  \[
    \boxed{S^p(r^*)=B}.
  \]

  \begin{itemize}
    \setlength{\itemsep}{0.52em}
    \item Los bonos comprados por los hogares deben ser iguales a los bonos emitidos por el gobierno.
    \item No hay inversión ni activos externos.
    \item Por tanto, el ahorro nacional es cero:
      \[
        S=S^p+S^g=S^p-B=0.
      \]
  \end{itemize}
  \notafuentes{Williamson (2018), ecuación 9-20.}
\end{frame}

\begin{frame}{El mercado de bienes también se vacía}
  Por definición,
  \[
    S^p=Y-T-C,
    \qquad
    B=G-T.
  \]
  Sustituyendo en \(S^p=B\),
  \[
    Y-T-C=G-T,
  \]
  por lo que
  \[
    \boxed{Y=C+G}.
  \]

  \begin{block}{Ley de Walras}
    Una vez satisfechas las restricciones presupuestarias, el vaciado del mercado de crédito implica el vaciado del mercado de bienes.
  \end{block}
  \notafuentes{Williamson (2018), ecuaciones 9-20--9-23.}
\end{frame}

% =============================================================================
\section{Equivalencia Ricardiana}
% =============================================================================

\begin{frame}{Un cambio en el momento de los impuestos}
  Considere una reducción de los impuestos actuales financiada con deuda:
  \[
    \Delta T<0,
    \qquad
    \Delta B=-\Delta T>0.
  \]
  Para satisfacer la restricción presupuestaria intertemporal, el gobierno tiene que recaudar en el periodo futuro la cantidad necesaria para pagar la deuda presente:
  \[
    \Delta T'=-(1+r)\Delta T>0.
  \]

  \begin{itemize}
    \item Para el sector privado, el ingreso disponible actual aumenta.
    \item Sin embargo, el ingreso disponible futuro disminuye en valor presente por la misma cantidad.
    \item La riqueza intertemporal del sector privado no cambia.
  \end{itemize}
  \notafuentes{Williamson (2018), pp. 337--339.}
\end{frame}

\begin{frame}{Teorema de equivalencia ricardiana}
  \begin{block}{Resultado}
    Manteniendo constante la trayectoria de gasto público, un cambio en la distribución temporal de impuestos no modifica el consumo ni la tasa de interés real de equilibrio.
  \end{block}

  Bajo una distribución uniforme de la carga tributaria,
  \[
    t+\frac{t'}{1+r}
    =\frac{1}{N}
    \left(G+\frac{G'}{1+r}\right).
  \]
  Por tanto, la restricción del hogar es
  \[
    c+\frac{c'}{1+r}
    =y+\frac{y'}{1+r}
    -\frac{1}{N}\left(G+\frac{G'}{1+r}\right),
  \]
  que no depende del momento en que se cobran los impuestos.
  \notafuentes{Williamson (2018), ecuaciones 9-24--9-26.}
\end{frame}

\begin{frame}[t]{La restricción del hogar no cambia}
  \begin{columns}[T,onlytextwidth]
    \begin{column}{0.41\textwidth}
      \small
      \begin{itemize}
        \item El recorte actual mueve la dotación de \(E_1\) a \(E_2\).
        \item El aumento futuro de impuestos compensa exactamente el recorte en valor presente.
        \item La restricción intertemporal y la elección \(A\) no cambian.
        \item El hogar ahorra íntegramente el recorte tributario.
      \end{itemize}
    \end{column}
    \begin{column}{0.55\textwidth}
      \centering
      \figura[5.15cm]{fig_9_16_equivalencia_ricardiana.png}
      \fuentefigura{Williamson (2018), figura 9.16.}
    \end{column}
  \end{columns}
  \notafuentes{Williamson (2018), figura 9.16 y pp. 338--339.}
\end{frame}

\begin{frame}{Ahorro privado y ahorro público}
  Aunque el consumo no cambia, la composición del ahorro sí:
  \[
    S^p=Y-T-C,
    \qquad
    S^g=T-G=-B.
  \]
  Ante un recorte tributario actual,
  \[
    \Delta S^p=-\Delta T>0,
    \qquad
    \Delta S^g=\Delta T<0.
  \]
  Por tanto,
  \[
    \Delta S=\Delta S^p+\Delta S^g=0.
  \]

  El mayor ahorro privado compensa exactamente la reducción del ahorro público.
  \notafuentes{Williamson (2018), pp. 338--340.}
\end{frame}

\begin{frame}[t]{Equivalencia ricardiana en el mercado de crédito}
  \begin{columns}[T,onlytextwidth]
    \begin{column}{0.40\textwidth}
      \small
      \begin{itemize}
        \item La deuda pública aumenta de \(B_1\) a \(B_2\).
        \item La oferta privada de crédito se desplaza exactamente en la misma magnitud.
        \item La tasa de interés permanece en \(r_1\).
        \item No existe desplazamiento del consumo privado.
      \end{itemize}
    \end{column}
    \begin{column}{0.56\textwidth}
      \centering
      \figura[4.95cm]{fig_9_17_mercado_credito.png}
      \fuentefigura{Williamson (2018), figura 9.17.}
    \end{column}
  \end{columns}
  \notafuentes{Williamson (2018), figura 9.17 y pp. 339--341.}
\end{frame}

\begin{frame}{¿Cuándo puede fallar la equivalencia ricardiana?}
  \begin{enumerate}
    \setlength{\itemsep}{0.52em}
    \item \textbf{Redistribución:} el recorte y los impuestos futuros recaen sobre hogares distintos.
    \item \textbf{Generaciones distintas:} quienes reciben el beneficio actual no pagan la deuda futura.
    \item \textbf{Impuestos distorsionantes:} cambiar su momento modifica precios marginales y decisiones.
    \item \textbf{Mercados de crédito imperfectos:} algunos hogares no pueden endeudarse o enfrentan tasas diferentes.
  \end{enumerate}

  \begin{alertblock}{Lectura correcta del teorema}
    Es un punto de referencia para identificar qué fricción hace que la deuda y el momento de los impuestos tengan efectos reales.
  \end{alertblock}
  \notafuentes{Williamson (2018), pp. 341--343.}
\end{frame}

\begin{frame}{La deuda pública no es riqueza neta bajo equivalencia}

  \begin{itemize}
    \setlength{\itemsep}{0.55em}
    \item La deuda pública es simultáneamente un activo privado en el presente y un pasivo fiscal futuro.
    \item Emitir deuda no crea recursos: modifica el momento de la tributación.
    \item La deuda sí puede tener efectos reales cuando redistribuye recursos, relaja restricciones financieras o altera distorsiones tributarias.
  \end{itemize}

  \centering
  \cidehighlight{Un recorte de impuestos financiado con deuda no es gratuito.}
  \notafuentes{Williamson (2018), pp. 340--343.}
\end{frame}

% =============================================================================
\section{Síntesis}
% =============================================================================

\begin{frame}{Resultados centrales}
  \begin{itemize}
    \setlength{\itemsep}{0.52em}
    \item La ecuación de Euler caracteriza la asignación de consumo intertemporal óptima del hogar.
    \item El ahorro permite suavizar consumo frente a variaciones del ingreso.
    \item El consumo depende de la riqueza intertemporal, no solo del ingreso corriente (Keynes).
    \item Los cambios permanentes del ingreso afectan más al consumo que los transitorios.
    \item La respuesta a la tasa real combina efectos sustitución e ingreso y dependen de la posición financiera del hogar.
    \item Introdujimos un nuevo precio, $r$, el cual se ajusta endógenamente para vaciar el mercado de crédito.
    \item Bajo condiciones estrictas, el momento de los impuestos es neutral (Equivalencia Ricardiana).
  \end{itemize}
\end{frame}

\begin{frame}{Referencias}
  \begin{itemize}
    \setlength{\itemsep}{0.8em}
    \item Williamson, Stephen D. (2018). \textit{Macroeconomics}, 6th ed. Pearson. Capítulo 9: ``A Two-Period Model: The Consumption-Savings Decision and Credit Markets''.
    \item Kurlat, Pablo (2020). \textit{A Course in Modern Macroeconomics}. Capítulo 6: ``Consumption and Saving''.
    \item Friedman, Milton (1957). \textit{A Theory of the Consumption Function}. Princeton University Press.
  \end{itemize}
\end{frame}





\appendix


\begin{frame}{Intuición del valor presente}
\hypertarget{app:VP}{} 
\small
  \textbf{El problema:} no se puede sumar \(y + y'\) directamente. Un bien del periodo actual y un bien del periodo futuro son bienes \textit{distintos}, así como no sumarías manzanas con naranjas o dolares con euros.

  \vspace{0.4em}
  \textbf{La solución:} convertir todo a una misma unidad usando el precio relativo entre ambos bienes, \(\frac{1}{1+r}\), y luego sí sumar.

  \vspace{0.5em}
  \begin{itemize}
    \item Si prestas un bien de hoy, el mercado te da \((1+r)\) bienes del futuro a cambio.
    \item Para obtener un bien del futuro, necesitas ahorrar \(\frac{1}{1+r}\) bienes de hoy.
    \item Por lo tanto, un bien del futuro vale \(\frac{1}{1+r}\) bienes de hoy.
    \item El valor presente \(y + \frac{1}{1+r}y'\) es la cantidad total de bienes en el tiempo a \textbf{valor de hoy}.
  \end{itemize}
  \vspace{0.4em}
  \cidehighlight{\(\frac{1}{1+r}\) es un \textbf{precio}: nos dice la tasa a la que el mercado está dispuesto a intercambiar un bien por otro.}


  \hyperlink{RP}{\beamerbutton{Regresar a Restricción Presupuestaria Intertemporal}}
\end{frame}

\begin{frame}{Intuición del valor presente: Un ejemplo numérico}
  Supongamos \(r = 10\%\).

  \vspace{0.4em}
  \begin{itemize}
    \setlength{\itemsep}{0.6em}
    \item Si tienes \$100 hoy y los prestas, mañana tienes
      \[
        \$100 \times (1+r) = \$100 \times 1.1 = \$110.
      \]
      Es decir: \textbf{\$100 hoy son \$110 mañana.}
    \item Al revés: ¿cuánto valen hoy \$110 que voy a recibir mañana? Es la
      pregunta inversa, así que dividimos entre \((1+r)\):
      \[
        \frac{\$110}{1+r} = \frac{\$110}{1.1} = \$100.
      \]
      Es decir: \textbf{\$110 mañana son \$100 hoy.}
  \end{itemize}

  \notafuentes{Kurlat (2020), sección 6.2.}
  \hyperlink{RP}{\beamerbutton{Regresar a Restricción Presupuestaria Intertemporal}}
\end{frame}

\begin{frame}{Intuición del valor presente: La idea general}
    \vspace{0.5em}
  \begin{block}{La idea general}
    \[
      \underbrace{y}_{\text{hoy}} \;\times\; (1+r) \;=\; \underbrace{y'}_{\text{mañana}}
      \qquad\Longleftrightarrow\qquad
      \underbrace{y'}_{\text{mañana}} \;\times\; \frac{1}{1+r} \;=\; \underbrace{y}_{\text{hoy}}
    \]
  \end{block}

  \cidehighlight{Multiplicar por \((1+r)\) lleva pesos de hoy a pesos de
  mañana; multiplicar por \(\dfrac{1}{1+r}\) trae pesos de mañana a pesos de
  hoy. Por eso \(\dfrac{1}{1+r}\) es el \textbf{precio} del consumo futuro.}
  \hyperlink{RP}{\beamerbutton{Regresar a Restricción Presupuestaria Intertemporal}}
\end{frame}


\begin{frame}{De dos periodos a muchos periodos}
  En un horizonte de \(T+1\) periodos, el hogar resuelve
  \[
    \max_{\{c_t,a_{t+1}\}_{t=0}^{T}}
    \sum_{t=0}^{T}\beta^t u(c_t)
  \]
  sujeto a
  \[
    a_{t+1}=(1+r_t)a_t+y_t-t_t-c_t,
    \qquad a_{T+1}=0.
  \]
  Para cada par de periodos consecutivos, la condición de optimalidad intertemporal es
  \[
    \boxed{u'(c_t)=\beta(1+r_{t+1})u'(c_{t+1})}.
  \]

  \begin{itemize}
    \item La ecuación de Euler determina la trayectoria completa del consumo.
    \item Al igual que en el modelo con dos periodos, la decisión óptima de consumo no depende de la trayectoria específica del ingreso.
  \end{itemize}
  \notafuentes{Kurlat (2020), sec. 6.3.}
\end{frame}

\begin{frame}{Horizonte infinito}
  Cuando \(T \to \infty\), el hogar resuelve
  \[
    \max_{\{c_t,a_{t+1}\}_{t=0}^{\infty}}
    \sum_{t=0}^{\infty}\beta^t u(c_t)
  \]
  sujeto a
  \[
    a_{t+1}=(1+r_t)a_t+y_t-t_t-c_t,
    \qquad a_0 \text{ dado},
  \]
  \[
    \lim_{T\to\infty}\;
    \frac{a_{T+1}}{\displaystyle\prod_{s=0}^{T}(1+r_s)} \;\geq\; 0
    \qquad \text{(condición de no-Ponzi)}.
  \]

  \vspace{0.3em}
  Para cada par de periodos consecutivos, la condición de optimalidad
  intertemporal es
  \[
    \boxed{u'(c_t)=\beta(1+r_{t+1})u'(c_{t+1})}.
  \]
\end{frame}

\end{document}
